\section{The closure model}

We treat the factory as a bill of materials and define its closure $C$ as the mass
fraction of its own parts it can manufacture on site. The complement, $1-C$, is the
mass fraction that must be imported as vitamins for every kilogram the factory builds.
Mass balance makes this exact rather than approximate: to add one kilogram of
installed capacity, the factory must be supplied with $1-C$ kilograms it cannot make.

The literature does not place $C$ at one. The canonical seed-factory study estimated
that \mbox{90-96\%} mass closure is achievable and argued that chasing the last few
percent to full closure is not worthwhile \cite{nasa-cp-2255-1980}. A recent revival
of the concept is more conservative, taking a near-term-realistic closure of about
70\% with chips, solar cells, and circuitry staying Earth-sourced
\cite{guided-self-replicating-factory-2021}. The modern synthesis agrees on the
principle: because the marginal part gets disproportionately harder to make as closure
rises, the sensible target is high-but-incomplete closure, not one
\cite{freitas-merkle-2004}. Across these designs the conclusion is the same. Imported
electronics are a permanent design feature, not a temporary compromise on the way to
full autonomy. The interesting question is therefore not \emph{whether} a factory
imports electronics but \emph{why} that particular part is the one it cannot make, and
what the resulting cap on closure costs.
